\appendix
\section{Setup details}
\label{app:setup}
\textbf{Outcome models.} All outcomes are $Y_i(z)=b_i+(\text{treatment part})+\varepsilon_i$ with a \emph{fixed} baseline per world, $b_i=u_i+0.5\,\tilde\ell_i+\eta_i$, where $u$ is a graph-smooth component (white noise passed through $W$ three times, rescaled to SD $\alpha_{\rm graph}$, default 1), $\tilde\ell_i$ the standardised log-degree, $\eta_i\sim N(0,0.5^2)$, and $\varepsilon_i\sim N(0,\sigma^2)$, $\sigma=1$. The mechanisms are:
\emph{shape}: $\tau z_i+\gamma h(f_i)$ with $h$ linear, concave $h(f)=\frac{1-e^{-kf}}{1-e^{-k}}$, convex $h(f)=\frac{e^{kf}-1}{e^{k}-1}$ ($k=3$ unless varied; $k\in[0.5,8]$ in the noiseless grid), or sigmoid (logistic with slope 10, rescaled to $h(0)=0,h(1)=1$);
\emph{count}: $\tau z_i+\gamma\,(Az)_i/\bar d$; \emph{twohop}: $\tau z_i+\gamma f_i+\gamma_2 f^{(2)}_i$ with $\gamma_2=0.5$ and $f^{(2)}$ the treated fraction among two-hop neighbours; \emph{hetero}: the direct effect $\tau(1+0.5\tilde\ell_i)$ with linear spillover; \emph{propagation}: $Y=(I-\rho W)^{-1}(b+\tau z)+\varepsilon$ for $\rho\in\{0.3,0.6,0.9\}$. Defaults are $\tau=1$, $\gamma=2$ (so $\tau_G=3$ for the shape mechanisms).
The GATE is the exact difference of the noise-free mean outcomes at $z=\mathbf 1$ and $z=\mathbf 0$ of the world.

\textbf{Estimators.} All OLS fits use the pseudo-inverse and predict at $z=\mathbf 1$ and $z=\mathbf 0$ with the same feature map. With $f_1$ the neighbour-treated fraction, $f_2$ the two-hop fraction, $c$ the neighbour count divided by the mean degree and $\ell$ the centred log-degree, the columns are: OLS$(z,f)$: $(1,z,f_1)$; OLS$_{\rm zf}$ adds $zf_1$; OLS$_{\rm cnt}$: $(1,z,c)$; OLS$_{\rm fd}$: $(1,z,f_1,\ell,z\ell,f_1\ell)$; OLS$_{f2}$: $(1,z,f_1,f_2)$; OLS quad: $(1,z,f_1,f_1^2,zf_1)$; OLS cubic adds $f_1^3$; OLS big: $(1,z,f_1,f_2,c,\ell,z\ell,f_1\ell,zf_1,f_1^2,f_2^2)$. ORACLE uses exactly the exposure columns of the mechanism ($(1,z,h(f_1))$ for shape, $(1,z,c)$ for count, $(1,z,f_1,f_2)$ for two-hop, $(1,z,z\ell,f_1)$ for heterogeneous direct effects); ORACLE$_D$ adds $\ell$, needed because baselines depend on degree; no oracle exists for propagation. GBM is scikit-learn histogram gradient boosting (150 iterations, depth 4, learning rate 0.1) on $(z,f_1,f_2,\log d,\text{count})$ with the same predict-at-corners contrast. HajekFull compares units that are treated with all neighbours treated against units that are control with all neighbours control, with exact design probabilities; HajekKappa ($\kappa=0.75$, Bernoulli only) uses ``treated with at least $\kappa$ of neighbours treated'' against ``control with at most $1-\kappa$ treated''. Both are undefined when no unit has the exposure; those cells are reported as unestimable (counts in the summaries).

\textbf{Variance estimators.} HC0 sandwich for the contrast; cluster-robust over fixed analysis clusters (Louvain partitions at resolution 0.5, 1, 2 or 5; 1 is the default); Conley-type graph-HAC with a uniform kernel over graph distance $\le1$ or $\le2$ (variance clipped at zero), computed only when $n\bar d^{2}<4\cdot10^{7}$, hence not for Stanford3's two-hop kernel.

\textbf{Grids and graphs.}
\begin{table}[h]
\caption{Grids: number of cells and replicates per cell (\texttt{results/tables/t13\_grid\_spec.csv}). $^\dagger$Designed after the first results were read (Appendix~\ref{app:post2}).}
\label{tab:grids}
\centering\scriptsize\setlength\tabcolsep{3pt}
\input{tab1_grids.tex}
\end{table}
\begin{table}[h]
\caption{Difference in means, OLS$(z,f)$ and the correctly specified regression by graph (concave spillover, $\gamma=2$, Bernoulli(0.5); \texttt{results/summary/e2.csv}).}
\label{tab:graphs}
\centering\scriptsize\setlength\tabcolsep{3pt}
\input{tab3_graphs.tex}
\end{table}
Stanford3: socfb-Stanford3, connected as downloaded ($n=11{,}586$, 568{,}309 edges, mean degree 98.1); ego-Facebook: largest connected component, $n=4039$, 88{,}234 edges, mean degree 43.7 (\texttt{results/tables/t13b\_real\_graphs.csv}). For real graphs only baselines, noise and assignments vary between worlds.

\begin{table}[h]
\caption{Closed form (\ref{eq:bias}) versus simulation, 144 noiseless cells; \texttt{results/summary/theory\_vs\_sim\_e9.csv}.}
\label{tab:theory}
\centering\scriptsize\setlength\tabcolsep{3pt}
\input{tab4_theory.tex}
\end{table}
\begin{table}[h]
\caption{Log--log fits of the SD on $n$ and mean degree $d$ (ER graphs, Bernoulli(0.5)); \texttt{results/summary/e3.csv}.}
\label{tab:scaling}
\centering\scriptsize\setlength\tabcolsep{3pt}
\input{tab5_scaling.tex}
\end{table}

\section{Derivations}
\label{app:theory}
\textbf{Bias formula (\ref{eq:bias}).} Take Bernoulli$(p)$ assignment, no self-loops, and noiseless outcomes $Y_i=b_i+\tau z_i+\gamma h(f_i)$. For unit $i$, $z_i$ is independent of $f_i=\mathrm{Bin}(d_i,p)/d_i$, and $\mathbb E f_i=p$ for all $i$. Pool units and assignments. Then $\mathrm{Cov}(z,f)=0$; $\mathrm{Cov}(b,\cdot)=0$ because $b_i$ is fixed and $f_i-p$ has mean zero given $i$; hence the least-squares coefficients on $(z,f)$ are $\tau$ for $z$ and $s=\mathrm{Cov}(h(f),f)/\mathrm{Var}(f)$ for $f$, with moments over the mixture of degrees. The predicted contrast is $\hat\tau\to\tau+s\gamma\,[f(\mathbf 1)-f(\mathbf 0)]=\tau+\gamma s$, since $f(\mathbf 1)-f(\mathbf 0)=1$, while $\tau_G=\tau+\gamma(h(1)-h(0))=\tau+\gamma$. This gives $\gamma(s-1)$. As $d_i\to\infty$, $f_i\to p$ and the delta method gives $\mathrm{Cov}(h(f),f)\approx h'(p)\mathrm{Var}(f)$, i.e.\ $s\to h'(p)$; for finite degrees $s$ is computed exactly from the binomial mixture (\texttt{code/scripts/theory\_bias.py}). For $h$ linear, $s=1$ and the bias is zero for any degree distribution. The sampling error of the OLS fit is an additional $O(n^{-1/2})$ term, which is why the formula is tested on noiseless cells with the finite-sample deviation reported alongside.

\textbf{Variance heuristic.} The contrast vector is $c=(0,1,1)$. With homoskedastic independent errors of variance $\sigma_r^2$ and $z$, $f$ uncorrelated, $\mathrm{Var}(\hat\tau)\approx\sigma_r^2\,c^\top(X^\top X)^{-1}c\approx\frac{\sigma_r^2}{n}\big[\frac{1}{p(1-p)}+\frac1{\mathrm{Var}(f)}\big]$, with $\mathrm{Var}(f)=p(1-p)\,\overline{1/d}$, so $\mathrm{Var}(\hat\tau)\approx\frac{\sigma_r^2}{n\,p(1-p)}(1+d_h)$, $d_h=1/\overline{1/d}$ the harmonic mean degree. We test the scaling in $n$, $d$, $p$ only, because neighbouring units share baselines and exposures, which violates the independence assumption and changes the level.

\section{Post-freeze grids e11--e13}
\label{app:post}
The three grids were run after the first results were seen, on a request from the planning step of the pipeline; the first-freeze summaries are unchanged. Rules are in \texttt{code/scripts/post\_freeze.py}, written (as reported by the simulation step; we could not verify the order from the commit history) before any cell of these grids was read; results are in \texttt{results/summary/e11.csv}, \texttt{e12.csv} and \texttt{results/tables/t14\_gamma\_dependence.csv}, \texttt{t15\_hajek.csv}.

\textbf{e11: bias against interference strength (H3, 64 cells, 400 estimates each).} SBM, BA, Stanford3 and ego-Facebook $\times$ concave and sigmoid spillover $\times$ $p\in\{0.1,0.5\}$ $\times$ $\gamma\in\{0.5,1,2,4\}$, noisy. The baseline and the noise do not depend on $\gamma$ and OLS is linear in the outcome, so the expected bias is exactly proportional to $\gamma$; the rule fits $\text{bias}=a+b\gamma$ per group and calls it proportional if $R^2\ge0.99$ and $|a|\le3\max(\mathrm{SE}_a,0.02)$. For OLS$(z,f)$, 14 of 16 groups are proportional (H3g supported). The slope per unit $\gamma$ is $-0.25$ (SBM, concave, $p=0.5$) and $+1.26$ ($p=0.1$), $+1.01$ and $-0.71$ for the sigmoid; BA $-0.23$, $+1.20$, $+0.89$, $-0.68$; Stanford3 $-0.08$, $+0.51$, $+0.29$, $-0.22$; ego-Facebook $-0.12$, $+0.88$, $+0.48$, $-0.40$ (slope SE 0.005--0.037). The sign and size are set by the shape and $p$ as in Sec.~\ref{sec:res-theory}; only the scale is $\gamma$. The two failing groups are both ego-Facebook (concave, $p=0.5$: $R^2=0.96$ with a small bias; sigmoid, $p=0.1$: intercept $+0.20\pm0.06$). The correctly specified regression with log-degree has $|\text{bias}|\le0.17$ at $\gamma=4$ in all groups, and the difference in means has slope $-1.00$. \emph{Negative result:} the quadratic regression OLS$(z,f,f^2)$ at $p=0.1$ has a $\gamma$-independent bias of $-9$ to $-10$ (SBM, BA), $-6$ (ego-Facebook) and $-1.9$ (Stanford3), an extrapolation instability (cf.\ e10), not a spillover effect.

\textbf{e12: Hájek baselines under correct and wrong exposure mappings (H6, 108 cells).} ER and WS (mean degree 4, 8) and BA ($m=2,4$), $n=3000$, Bernoulli(0.5), Bernoulli(0.8) and Louvain clusters; mechanisms: threshold (step) with $\kappa=1$ and $\kappa=0.75$ (the full-neighbourhood mapping is correct for both, the $\kappa=0.75$ mapping only for the second), linear and sigmoid one-hop shapes, two-hop, and propagation ($\rho=0.6$); noisy; $\tau_G=3$ for step and shape mechanisms. Both estimators are estimable in 89\% (full) and 92\% ($\kappa$) of cells ($n_{\text{rep}}\ge200$ of 400); all other cells are Bernoulli(0.8), where fully control-exposed units are rare. $\kappa$-Hájek is undefined under cluster designs. \emph{H6c (correct mapping):} unbiased ($|\text{bias}|\le\max(3\,\mathrm{SE},0.1)$) in 31 of 43 estimable cells (72\%; rule $\ge75\%$): unclear. By design (post hoc, same cells): Bernoulli(0.5) 18 of 18, clusters 12 of 12, Bernoulli(0.8) 1 of 13, where the ratio estimator is biased upward ($+0.5$ to $+0.7$ full, $+0.28$ $\kappa$) because of the few control-exposed units. \emph{H6d (wrong mapping):} biased ($>\max(3\,\mathrm{SE},0.2)$) in 44 of 65 cells (68\%): unclear; post hoc Bernoulli(0.5) 24 of 30, Bernoulli(0.8) 12 of 23, clusters 8 of 12. At Bernoulli(0.5) the mean bias is $-0.88$ (full on two-hop), $-0.65$ (full on propagation), $-0.57$ ($\kappa$ on linear) and $-1.38$ ($\kappa$ on the $\kappa=1$ step), but $-0.13$ for $\kappa$ on the sigmoid, because $f\ge0.75$ already lies on its saturated part. The full-neighbourhood estimator is not wrong on one-hop shapes (bias $0.03$ linear, $-0.01$ sigmoid): its full-exposure contrast equals the GATE for any shape. \emph{H6e (descriptive):} Hájek has lower RMSE than OLS$(z,f)$ in 32\% of 96 (full) and 30\% of 66 ($\kappa$) estimable cells (median RMSE ratio 1.10 and 1.28). The correctly specified regression with log-degree is unbiased ($|\text{bias}|\le0.04$) in the step, shape and two-hop mechanisms.

\begin{figure}[tbp]
\centering\includegraphics[width=\linewidth]{fig9_gamma_dependence.pdf}
\caption{Bias against interference strength $\gamma$ under concave and sigmoid spillover (post-freeze grid e11; marker = graph). The difference in means has bias $-\gamma$; OLS$(z,f)$ has bias proportional to $\gamma$ whose sign and size are set by shape and $p$ (14 of 16 groups proportional; both exceptions are ego-Facebook); the correctly specified regression is unbiased. OLS quad is omitted (Appendix~\ref{app:post}).}
\label{fig:gamma}
\end{figure}
\begin{figure}[tbp]
\centering\includegraphics[width=\linewidth]{fig10_hajek.pdf}
\caption{Bias relative to the GATE of the exposure-mapping (Hájek) estimators, the difference in means, OLS$(z,f)$ and the correctly specified regression on low-degree graphs (post-freeze grid e12), by mechanism and design (marker = graph). Filled: the exposure mapping or model is correct for the mechanism; open: wrong.}
\label{fig:hajekfig}
\end{figure}

\textbf{e13: second seed (217 cells).} The cells of the grids that carry H2, H4 and H5 were rerun with a new seed (\texttt{results/tables/t16\_second\_seed.csv}; the replication rule, at least 75\% of matched Stanford3 cells with Louvain RMSE above random-cluster RMSE, was fixed before any e13 cell was read, according to the simulation step). H5: Louvain RMSE exceeds random-cluster RMSE on Stanford3 in 10/12 matched cells in the first seed and 11/12 in the second, and in 0/12 on ego-Facebook in both (verdict ``replicated''); the cell-mean RMSEs flip between seeds (Louvain 0.92 vs 1.11 first seed, 0.99 vs 0.84 second), dominated by the 10-cluster resolution. H2: correlation of simulation and closed form $1.0000$, maximum deviation $0.058$, all 144 cells within 3 SE. H4: ER exponents $-0.506\pm0.008$ in $n$ and $+0.508\pm0.009$ in $d$.

\section{Post-freeze analyses e14--e21}
\label{app:post2}
These eight grids (427 cells) and their analyses were designed after the first results and after e1--e13 had been read, to answer questions raised on the first version. They have \emph{no} decision rule or threshold: each item is an estimate with its Monte-Carlo uncertainty. Summaries are \texttt{results/summary/e14.csv}--\texttt{e21.csv}, analysis tables \texttt{results/tables/t17}--\texttt{t25}, and per-replicate data for every grid are in \texttt{results/compact/} (Appendix~\ref{app:repro}).

\textbf{Variance law across graphs (e21 and \texttt{t17}).} Table~\ref{tab:vlaw}.  See Sec.~\ref{sec:res-theory}. With independent errors only ($\gamma=0$, no baseline), observed/predicted SD with $d_h$ is between 0.93 and 1.04 on all six graphs (max/min 1.12); with the arithmetic mean degree, max/min is 2.82 (Stanford3: 0.185 predicted, 0.064 observed). With the default outcome, observed/predicted is 1.29 (ER), 1.44 (BA), 2.40 (SBM), 2.60 (Stanford3), 2.59 (WS), 5.1 (ego-Facebook). The graph-smooth baseline multiplies the SD of OLS$(z,f)$ by 7.2 (ego-Facebook), 3.9 (WS), 3.7 (SBM), 3.6 (Stanford3), 2.1 (BA), 1.8 (ER). Stanford3 and ER have the same SD ($\approx0.26$) because a lower $d_h$ (10.8 against 19) is cancelled by a larger baseline effect.

\textbf{Stanford3 cluster-design mechanism (\texttt{t18}).} Table~\ref{tab:expo}.  Within-arm variance of $f$ under Louvain clusters of resolutions 1--20: Stanford3 0.021--0.028 (Bernoulli 0.023, random clusters of the same size 0.020--0.023) although $\mathrm{corr}(z,f)=0.70$--$0.88$; SBM 0.0046 (Bernoulli 0.013, $\mathrm{corr}=0.985$); ego-Facebook 0.0035--0.040 (Bernoulli 0.0225).

\textbf{Matched degree-distribution test (e14, \texttt{t23}).} Table~\ref{tab:heavy}.  28 cells, $n=5000$, mean degree 20, Bernoulli(0.5), 8 worlds $\times$ 50 assignments, linear and concave truth, $\gamma=2$. SD of OLS$(z,f)$ with linear truth and $\gamma=2$: BA ($m=10$) 0.255, configuration model on the BA degree sequence 0.248, ER 0.276, 20-regular 0.240, power laws with exponents 2.2, 2.5, 3.0 and degrees $\le500$: 0.144, 0.164, 0.175 (at $\gamma=0$: 0.259, 0.243, 0.263, 0.236, 0.147, 0.157, 0.163); $|\text{bias}|\le0.014$ for linear truth. Concave bias: BA $-0.452\pm0.012$, configuration model $-0.487\pm0.018$. The power laws also have a larger share of low-degree nodes ($d_h$ 3.6--4.5, clustering 0.05--0.08).

\textbf{Coverage (e15, \texttt{t19}, \texttt{t19b}).} Table~\ref{tab:cov2}.  Same 24 cells (6 graphs $\times$ linear/concave $\times$ 2 designs), $4\times150$ estimates, as e6 plus the block bootstrap and the Stanford3 2-hop HAC. Mean coverage of the GATE / of $E[\hat\tau]$ over the 12 cells of a design (Bernoulli(0.5); Louvain clusters): HC0 0.54/0.66; 0.36/0.37. HAC-1 0.66/0.78; 0.56/0.57. HAC-2 0.82/0.92; 0.73/0.74. Cluster-robust 0.75/0.86; 0.80/0.82. Block bootstrap (SD) 0.76/0.87; 0.83/0.85, (percentile) 0.75/0.86; 0.82/0.83. Minimum over cells of the GATE coverage: 0.31 (HC0, cluster-robust, bootstrap, Bernoulli).

On Stanford3 (concave/linear truth) the coverage of the GATE under Bernoulli is 0.93/0.95 (HAC-2), 0.84/0.87 (cluster-robust), 0.82/0.83 (HC0); under Louvain resolution 1, 0.57/0.59 (HAC-2), 0.73/0.76 (cluster-robust) and 0.84--0.92 (bootstrap).

\textbf{Estimator scope (e16, \texttt{t25}).} Table~\ref{tab:scope}.  The 48 Bernoulli cells of e2 rerun with the same seeds (OLS$(z,f)$, OLS$(z,f,zf)$ and OLS$_\text{big}$ reproduce e2 to the last digit). RMSE relative to OLS$(z,f)$, median over cells: OLS$(z,f,zf)$ 1.03 (1.00 at $p=0.5$, 1.09--1.72 at $p=0.1$); ridge 6.3 and lasso 6.5, against 6.6 for unpenalised OLS$_\text{big}$, and never below 1; boosting with depth $\{2,3,6\}$, learning rate $\{0.03,0.1,0.3\}$ and $\{100,300\}$ iterations chosen by cross-validation on the first replicate of each world: 2.03 at Bernoulli(0.5) (24 cells), against 1.71 for the untuned boosting over 48 cells; it wins only where OLS$(z,f)$ is badly biased (concave at $p=0.1$: 0.51; sigmoid: 0.40--0.59). Median $|\text{bias}|$ of tuned boosting 0.44 against 0.82 untuned.

\textbf{Crossover and variance ratio (e17, \texttt{t21}).} 135 cells: SBM, Bernoulli(0.5), linear truth, $n\in\{10^3,5\cdot10^3,2\cdot10^4\}$, $d\in\{10,20,40\}$, $\sigma\in\{0.5,1,2\}$, $\gamma\in\{0,0.25,0.5,1,2\}$. The bias slope of the difference in means is $-1.006$ to $-0.995$ in all 27 groups. $\gamma^*$ is the model-based crossover of RMSE (validated against interpolation of observed RMSEs); its values above the $\gamma$ grid (maximum 2.2 at $n=10^3$, $d=40$, $\sigma=2$) are extrapolations. Log--log fits: $\gamma^*\propto n^{-0.49\pm0.01}d^{1.01\pm0.02}\sigma^{0.09\pm0.02}$ ($R^2=0.997$), SD ratio at $\gamma=0$ $\propto n^{0.01\pm0.01}d^{1.01\pm0.02}\sigma^{-0.31\pm0.02}$ ($R^2=0.989$; 4.0 to 26.5, median 11.1).

\textbf{Variance decomposition (e20, \texttt{t22}).} SBM full factorial: $n\{10^3,5\cdot10^3,2\cdot10^4\}\times d\{10,20,40\}\times p\{0.1,0.5\}\times\sigma\{0.5,2\}\times$ graph-smooth baseline SD $\{0,2\}\times$ truth $\{$linear, concave$\}$, $\gamma=2$. First-order/total-effect shares (Sobol indices of cell means; MC uncertainty $\pm0.01$) for log RMSE of OLS$(z,f)$: baseline $0.22/0.33$, shape $0.20/0.36$, $p$ $0.16/0.20$, $n$ $0.13/0.17$, $d$ $0.09/0.12$, $\sigma$ $0.01/0.04$; the largest interaction is baseline $\times$ shape (0.06). For the SD: baseline $0.46/0.51$, $n$ $0.27/0.27$, $d$ $0.14/0.16$, $p$ $0.05/0.05$, $\sigma$ $0.03/0.06$, shape $0.00/0.00$. Two levels for four of the six factors; the shares describe the chosen levels.

\textbf{Uncertainty with four worlds (e18, \texttt{t20}).} Table~\ref{tab:worlds}.  12 cells (six graphs $\times$ linear/concave, Bernoulli(0.5)) with 20 worlds $\times$ 25 assignments (500 estimates). Intraclass correlation of the OLS$(z,f)$ error across worlds: median 0.000--0.004, maximum 0.035 (e18, concave) and 0.016 (e2); for the correctly specified regression up to 0.165 (e18) and 0.105 (e2, concave); the difference in means $\approx0$. Using subsets of $W$ worlds against the remaining worlds as reference, the interval bias$\pm1.96\,$SE with the aggregation rule (larger of replicate- and world-level SE) covers the reference in $0.968$ of draws at $W=4$, the world-level SE alone in $0.850$ ($0.940$ with a $t_3$ quantile; its 10th--90th percentile range is 0.014--0.079 around a median 0.031), the replicate-level SE alone in $0.958$; at $W=12$ the world-level SE covers $0.921$ ($0.946$ with $t$). This analysis was run for OLS$(z,f)$ only; with larger intraclass correlation (correctly specified model, concave) the world-level term matters more. Hierarchical-bootstrap 95\% intervals for the SD ratio of OLS$(z,f)$ to the difference in means at $\gamma=0$ (Bernoulli(0.5), $B=400$): SBM 10.1 [9.3, 11.1], BA 5.5 [4.8, 6.3], WS 12.3 [11.3, 13.6], ER 6.5 [5.8, 7.4], Stanford3 9.2 [7.9, 10.5], ego-Facebook 18.6 [16.9, 20.7].

\textbf{Design-based baselines under cluster designs at higher degree (e19, \texttt{t24}).} Louvain resolutions 1 and 5; SBM ($d=20$), ego-Facebook (44), Stanford3 (98); linear/concave, $\gamma=2$, 400 estimates. The full-exposure Hájek estimator is estimable in all but 2--6 of 400 replicates (ego-Facebook), with $|\text{bias}|\le0.05$ and RMSE relative to OLS$(z,f)$ of 1.19--1.25 (SBM), 1.32--1.47 (ego-Facebook), 0.89--0.96 (Stanford3, concave) and 1.03/2.01 (Stanford3, linear, resolutions 1/5); the correctly specified regression has the lowest RMSE in 11 of 12 cells (the difference in means in the twelfth). Fractional-exposure Hájek estimators ($\kappa=0.75,0.9$) are biased ($-0.04$ to $-0.77$) because they estimate a contrast at exposure $\kappa$, not 1. The first-version statement that design-based baselines are impractical at $d\ge20$ holds for Bernoulli designs only.

\begin{table}[h]
\caption{Coverage of nominal 95\% intervals of OLS$(z,f)$ for the GATE and for $E[\hat\tau]$ by interval type (e15): mean over 12 cells (six graphs $\times$ linear/concave, $\gamma=2$) with a bootstrap 95\% interval over cells, and min--max. Source: \texttt{t19b\_coverage\_by\_se\_type.csv}.}
\label{tab:cov2}
\centering\scriptsize\setlength\tabcolsep{3pt}
\resizebox{\linewidth}{!}{\input{tab6_coverage_se_types.tex}}
\end{table}
\begin{table}[h]
\caption{Observed SD of OLS$(z,f)$ against the variance law with harmonic-mean degree $d_h$ and mean degree $d$; $r$ = observed/predicted (e1, e3, e21). Source: \texttt{t17*.csv}.}
\label{tab:vlaw}
\centering\scriptsize\setlength\tabcolsep{3pt}
\resizebox{\linewidth}{!}{\input{tab7_variance_law.tex}}
\end{table}
\begin{table}[h]
\caption{Matched degree-distribution test (e14; $n=5000$, mean degree 19--20, Bernoulli(0.5)): SD and its ratio to the $d_h$ law at $\gamma=0$; bias (MC SE), SD and RMSE at $\gamma=2$. Source: \texttt{t23\_heavy\_tail.csv}.}
\label{tab:heavy}
\centering\scriptsize\setlength\tabcolsep{3pt}
\resizebox{\linewidth}{!}{\input{tab8_heavy_tail.tex}}
\end{table}
\begin{table}[h]
\caption{Estimator scope (e16): median RMSE relative to OLS$(z,f)$ and, after the slash, share of cells in which the estimator beats it, by spillover shape and $p$ (six graphs per entry); last column over all cells (number of cells). Ratios above 1 favour OLS$(z,f)$; ``--'': not run. Source: \texttt{t25\_estimator\_scope.csv}.}
\label{tab:scope}
\centering\scriptsize\setlength\tabcolsep{3pt}
\resizebox{\linewidth}{!}{\input{tab9_estimator_scope.tex}}
\end{table}
\begin{table}[h]
\caption{Variation of the neighbour-treated fraction $f$ under Bernoulli(0.5), Louvain-cluster and size-matched random-cluster designs (res.: resolution): variance of $f$, within-arm variance, $\mathrm{corr}(z,f)$ and SD of OLS$(z,f)$ in e4. Source: \texttt{t18\_exposure\_variation.csv}.}
\label{tab:expo}
\centering\scriptsize\setlength\tabcolsep{3pt}
\resizebox{\linewidth}{!}{\input{tab10_exposure_variation.tex}}
\end{table}
\begin{table}[h]
\caption{(a) Subsets of $W$ of the 20 worlds of e18 (12 cells): median SE of the bias by the aggregation rule, world-level only and replicate-level only, and how often bias $\pm1.96$ SE (or $\pm t$ SE) contains the 20-world mean. (b) Intraclass correlation of the error between worlds and design effect (DE) of the replicate-level SE, median and maximum over 12 cells per row. Sources: \texttt{t20b\_worlds\_needed.csv}, \texttt{t20\_variance\_components.csv}.}
\label{tab:worlds}
\centering\scriptsize\setlength\tabcolsep{3pt}
\resizebox{\linewidth}{!}{\input{tab11_worlds_icc.tex}}
\end{table}

\section{Hypotheses, decision rules and verdicts}
\label{app:hyp}
\textbf{When the rules were fixed.} The directional predictions H1--H8 were written in the study plan before any simulation was run. The numerical thresholds of the decision rules were written into \texttt{code/scripts/findings.py} \emph{after} the grids e1--e10 had completed, and were not changed to resolve ``unclear'' verdicts; the verdicts are therefore not pre-registered tests in the strict sense. The rules for the post-freeze grids e11 and e12 were written in \texttt{code/scripts/post\_freeze.py} before any of their cells was read (a statement we could not verify from the commit history, which does not order the two); those grids were nevertheless run after the first results were seen, at the request of a reviewing step of the pipeline, and are labelled as additions. Rules are those in \texttt{results/hypotheses.csv} and \texttt{results/hypotheses\_postfreeze.csv}. Part~A of the table lists consistency checks (identities of the model or checks of the code, which are expected to pass; ``passes'' says our implementation agrees with the algebra, not that the world behaves a certain way); parts~B and~C list empirical hypotheses, with verdicts ``supported'', ``refuted'' or ``unclear'' (the rule neither confirms nor refutes). Descriptive tests have no threshold.
{\scriptsize\begin{longtable}{@{}p{0.11\linewidth}p{0.33\linewidth}p{0.4\linewidth}p{0.08\linewidth}@{}}
\toprule ID & test & decision rule & verdict \\
\endhead
\multicolumn{4}{@{}l}{\emph{A. Consistency checks: identities of the model or of the code (a failure would indicate an error, a pass is not evidence about the world)}} \\ \midrule
H1 & noiseless ORACLE bias, e8 (all graphs/designs/shapes run) & supported if max$|$bias$|$$<$1e-8 (machine precision); refuted if any $|$bias$|$$>$5 MC SE & passes \\
H1 & noisy ORACLE\_D bias in units of its MC SE, e1,e2,e3,e5 & supported if share $|$z$|$$>$3 below 0.05 (expect $\sim$0.003 under H0) & passes \\
H2 & e9 noiseless OLS\_f bias vs exact-binomial closed form (theory\_bias.py) & supported if corr$>$0.99 and median$|$dev$|$$<$0.01 (bias scale 0.1-5); report share within 3 MC SE & passes \\
H3b & e9 noiseless $|$bias$|$ vs curvature (gamma fixed) & supported if median rho$>$0.8; gamma-dependence is NOT tested here (e9 fixes gamma) - see note & passes \\
H6a & e1 DM bias slope in gamma (truth: GATE = tau + gamma; DM sees tau only under Bernoulli) & supported if Bernoulli slopes within 0.05 of -1 (proportional to gamma) & passes \\
H3g & e11 OLS\_f: bias vs gamma in {0.5,1,2,4}, per graph x shape x design (post-freeze) & per group weighted line bias = a + b*gamma; proportional if R2 $\ge$ 0.99 and $|$a$|$ $\le$ 3 max(SE\_a, 0.02); supported if $\ge$ 75\% of groups with $|$bias(4)$|$$>$0.2 & passes \\
\midrule \multicolumn{4}{@{}l}{\emph{B. Empirical hypotheses (first freeze)}} \\ \midrule
H3a & e10 concave spillover: OLS\_f bias vs n (SBM, 1k-20k, Bernoulli) & supported if $|$bias$|$ at the largest n is $\ge$70\% of that at n=1000 for both p (sampling error would fall like n\^{}-1/2, to $\sim$22\%) & supported \\
H4a & e3 ER Bernoulli(0.5), OLS\_f: log SD $\sim$ log n + log d & supported if both exponents within 0.05 of prediction & supported \\
H4b & e3 SBM + Louvain-cluster design, OLS\_f: log SD $\sim$ log n + log d & descriptive (the prediction is stated for Bernoulli designs) & unclear \\
H4c & e4 SD vs p, OLS\_f, linear truth & supported if slope within [0.7,1.3] (pred 1) for every graph & unclear \\
H5a & e4 matched-resolution comparison of Louvain vs random clusters (OLS\_f RMSE; DM $|$bias$|$) & supported if good clusters have lower RMSE in $\ge$75\% of matched cells and random clusters show no gain over Bernoulli & unclear \\
H5b & e4 planted-block quality (within-block edge fraction 0.8/0.9/0.97), linear truth & descriptive: monotone improvement in RMSE with quality would support H5 & unclear \\
H6b & e1 DM wins RMSE at gamma=0, loses at larger gamma & supported if DM wins at gamma=0 in all combinations and OLS\_f wins by gamma$\le$4 & supported \\
H6c & Hajek / exposure-mapping baselines estimable? & reported as unestimable when n\_rep$<$10 (full-neighbourhood exposure units absent at this degree) & unclear \\
H7a & e6 mean 95\% CI coverage by SE type (degenerate cell e6 BA acl0.5 excluded: one analysis cluster) & supported if HC0 coverage $<$0.85 for correctly specified/linear cells and cluster-robust/HAC raise it; coverage loses under bias (concave OLS\_f $<$ 0.85) & supported \\
H8a & e2 graph family: OLS\_f RMSE (Bernoulli 0.5, linear truth) & supported if both heavy-tailed graphs (BA, Stanford3) have RMSE above the median of the six graphs; refuted if both are below it & refuted \\
H8b & e7 two-hop interference fitted with 1-hop features (OLS\_f) vs the 2-hop-aware ORACLE\_D, Bernoulli 0.5 & supported if OLS\_f $|$bias$|$$>$0.3 on all graphs & supported \\
C2a & e1 gamma=0, Bernoulli(0.5): SD of OLS\_f over SD of DM (no interference, both unbiased) & supported if the ratio exceeds 3 on every graph & supported \\
C2b & e2 OLS\_f: bias vs variance by spillover shape (gamma=2, 6 graphs x 4 designs) & bias dominates for sigmoid, not for linear (bias-dominated share 0 for linear) & supported \\
C2c & e2 RMSE comparisons over 96 graph x shape x design cells & descriptive (ranking of estimators by RMSE) & unclear \\
\midrule \multicolumn{4}{@{}l}{\emph{C. Empirical, post-freeze additions (run after the first results were seen)}} \\ \midrule
H6c & e12 Hajek bias under a CORRECT exposure mapping (step mechanism) & supported (unbiased) if $\ge$ 75\% of estimable (n\_rep $\ge$ 200) cells & unclear \\
H6d & e12 Hajek bias under a WRONG exposure mapping (HajekFull on two-hop/propagation; HajekKappa(0.75) on linear/sigmoid shape and on step kappa=1) & supported (biased) if $\ge$ 75\% of estimable cells & unclear \\
H6c-posthoc & H6c stratified by design (post-hoc reading of the same cells; the verdict above is the pre-stated rule) & descriptive & reported \\
H6d-posthoc & H6d stratified by design (post-hoc) & descriptive & reported \\
H6e & e12 RMSE of Hajek vs OLS\_f over estimable cells & descriptive (no verdict threshold) & reported \\
\bottomrule
\end{longtable}
}

\section{Reproduction and provenance}
\label{app:repro}
\textbf{Figure and table map} (paper number: file in \texttt{figures/out/}). Figs.~\ref{fig:strength}: \texttt{fig1\_strength}; \ref{fig:theory}: \texttt{fig2\_bias\_theory}; \ref{fig:bias-n}: \texttt{fig3\_bias\_vs\_n}; \ref{fig:var}: \texttt{fig4\_variance\_scaling}; \ref{fig:cov}: \texttt{fig6\_coverage}; \ref{fig:fail}: \texttt{fig7\_failure\_map\_relrmse\_sbm}; \ref{fig:rank}: \texttt{fig8\_factor\_ranking}; \ref{fig:cross}: \texttt{fig11\_crossover\_anova}; \ref{fig:gamma}: \texttt{fig9\_gamma\_dependence}; \ref{fig:hajekfig}: \texttt{fig10\_hajek}; \ref{fig:strength-all}: \texttt{fig1\_strength\_all}; \ref{fig:cluster}: \texttt{fig5\_cluster\_design}; Figs.~\ref{fig:fail-ba}--\ref{fig:fail-egofacebook}: \texttt{fig7\_failure\_map\_relrmse\_}\{\texttt{ba, er, ws, stanford3, egofacebook}\}. Tables~\ref{tab:est}: \texttt{tab2\_estimators}; \ref{tab:grids}: \texttt{tab1\_grids}; \ref{tab:graphs}: \texttt{tab3\_graphs}; \ref{tab:theory}: \texttt{tab4\_theory}; \ref{tab:scaling}: \texttt{tab5\_scaling}; \ref{tab:cov2}: \texttt{tab6\_coverage\_se\_types}; \ref{tab:vlaw}: \texttt{tab7\_variance\_law}; \ref{tab:heavy}: \texttt{tab8\_heavy\_tail}; \ref{tab:scope}: \texttt{tab9\_estimator\_scope}; \ref{tab:expo}: \texttt{tab10\_exposure\_variation}; \ref{tab:worlds}: \texttt{tab11\_worlds\_icc}. Each is produced by the script of the same name (tables: \texttt{make\_tables.py}) in \texttt{figures/scripts/} from \texttt{results/summary/}; \texttt{figures/MANIFEST.md} lists inputs. \textbf{Environment.} Python 3.11.2; numpy 2.4.6, scipy 1.17.1, pandas 3.0.6, scikit-learn 1.9.1, networkx 3.6.1 (\texttt{code/requirements.lock}); CPU only. \textbf{Data.} socfb-Stanford3 zip SHA-256 \texttt{b1221e41\dots8720d} and SNAP ego-Facebook SHA-256 \texttt{125e84db\dots035d7}; outcomes are simulated, no outcome data exist. \textbf{Per-replicate data.} \texttt{results/compact/} holds one row per cell, world, replicate and estimator for all 21 grids (about 70 MB); \texttt{aggregate.py --compact} rebuilds every per-cell summary from it, equal to the shipped summaries to a relative $3.5\cdot10^{-7}$, so all tables and figures can be regenerated without rerunning the simulations. \textbf{Seeds.} Each cell has a fixed integer seed; world $w$ uses \texttt{default\_rng(seed*1000+w)}. \textbf{Commands and runtime.} \texttt{REPRODUCE.md} lists the commands (run experiment, aggregate, tables, findings, factor ranking, figures); grids e1--e10 take about 58 CPU-hours (e9 1.4 h, e2 15 h, e7 11 h) and the post-freeze grids e11--e21 about 8 more (e13 3.3 h, e16 2.5 h, the others under 0.7 h each; \texttt{results/RUNTIME.csv}). The unit tests (16, 9 s) run with \texttt{pytest code/tests}. \textbf{Provenance of the text.} Every figure and table is made by a script listed in \texttt{figures/MANIFEST.md} from \texttt{results/summary}; the numbers quoted in the text are read from \texttt{results/} by \texttt{paper/check\_numbers.py}, which fails if a quoted value differs from the committed tables.

\textbf{Access to excluded material.} The task excludes the published paper on this exact question. No literature search on the question itself was run for this report; the bibliographic look-ups (arXiv and Crossref metadata of the starter-list works and standard references for graph generators and datasets) did not encounter it. The simulation step reports no encounter either (\texttt{REPRODUCE.md} \S8), and the literature-review step's brief states that the question was deliberately not searched and records no encounter.

\section{Additional figures}
\label{app:figs}
\begin{figure}[tbp]
\centering\includegraphics[width=\linewidth]{fig1_strength_all.pdf}
\caption{Bias and RMSE against interference strength $\gamma$ on all six graphs (linear truth, Bernoulli(0.5)).}
\label{fig:strength-all}
\end{figure}
\begin{figure}[tbp]
\centering\includegraphics[width=\linewidth]{fig8_factor_ranking.pdf}
\caption{One-factor ranges (not a ranking; see Sec.~\ref{sec:res-fail}). Range of the RMSE of OLS$(z,f)$ (filled) and of the correctly specified regression (open) over each one-factor sweep; ratio of max to min on the right.}
\label{fig:rank}
\end{figure}
\begin{figure}[tbp]
\centering\includegraphics[width=\linewidth]{fig5_cluster_design.pdf}
\caption{Graph (solid) versus size-matched random (dashed) clusters across resolution, concave spillover: share of units with $f\ge0.9$ (exposure support), bias, RMSE; dotted: Bernoulli(0.5).}
\label{fig:cluster}
\end{figure}
\begin{figure}[tbp]
\centering\includegraphics[width=\linewidth]{fig7_failure_map_relrmse_ba.pdf}
\caption{Relative RMSE failure map on BA (same layout as Fig.~\ref{fig:fail}).}
\label{fig:fail-ba}
\end{figure}
\begin{figure}[tbp]
\centering\includegraphics[width=\linewidth]{fig7_failure_map_relrmse_er.pdf}
\caption{Relative RMSE failure map on ER (same layout as Fig.~\ref{fig:fail}).}
\label{fig:fail-er}
\end{figure}
\begin{figure}[tbp]
\centering\includegraphics[width=\linewidth]{fig7_failure_map_relrmse_ws.pdf}
\caption{Relative RMSE failure map on WS (same layout as Fig.~\ref{fig:fail}).}
\label{fig:fail-ws}
\end{figure}
\begin{figure}[tbp]
\centering\includegraphics[width=\linewidth]{fig7_failure_map_relrmse_stanford3.pdf}
\caption{Relative RMSE failure map on Stanford3 (same layout as Fig.~\ref{fig:fail}).}
\label{fig:fail-stanford3}
\end{figure}
\begin{figure}[tbp]
\centering\includegraphics[width=\linewidth]{fig7_failure_map_relrmse_egofacebook.pdf}
\caption{Relative RMSE failure map on ego-Facebook (same layout as Fig.~\ref{fig:fail}).}
\label{fig:fail-egofacebook}
\end{figure}
